\frfilename{mt2a.tex}
\versiondate{3.8.15}

\def\RoverC{\hbox{\biggerscriptfonts${{\Bbb R}\atop{\Bbb C}}$}}

\def\chaptername{Lampiran}
\def\sectionname{Pendahuluan}

\gdef\topparagraph{}
\gdef\bottomparagraph{Lampiran Jilid 2 {\it pendahuluan}}

\centerline{\bf Lampiran Jilid 2}

\medskip

\centerline{\bf Fakta-Fakta Berguna}

\medskip

Dalam penulisan jilid ini, saya mendapati bahwa kita memerlukan cukup banyak
konsep dan fakta dari berbagai cabang matematika. Hampir semuanya merupakan
bagian dari teori-teori penting yang sudah mapan, yang memiliki banyak buku
ajar unggul dan yang sangat saya harapkan kelak akan Anda pelajari secara
mendalam. Meskipun demikian, saya enggan langsung mengirim Anda ke kuliah
topologi umum, analisis fungsional, dan teori himpunan, seolah-olah semua itu
merupakan prasyarat mutlak bagi pekerjaan kita di sini, di samping analisis
real dan aljabar linear dasar. Karena itu saya menulis Lampiran ini, yang
menyajikan hasil-hasil yang benar-benar kita perlukan pada suatu titik dalam
jilid ini. Sebagian besar hasil tersebut memang elementer---bahkan beberapa
begitu elementer sehingga uraian baku mengenai bidangnya tidak selalu
menuliskannya secara terperinci.

Walaupun saya tidak mengajukan buku ini sebagai tempat yang tepat untuk
mempelajari satu pun dari topik-topik tersebut, saya telah berusaha
menyajikannya dengan cara yang mudah Anda padukan dengan pendekatan-pendekatan
baku. Saya tidak mengharapkan siapa pun membaca karya ini secara sistematis,
dan saya berharap rujukan-rujukan dalam bab-bab utama jilid ini memadai untuk
mengarahkan Anda kepada butir khusus yang Anda perlukan.

\discrpage

